\documentclass[11pt,a4paper]{article}
\usepackage[T1]{fontenc}
\usepackage[utf8]{inputenc}
\usepackage{lmodern,microtype,amsmath,amsthm,amssymb,booktabs,fancyvrb}
\usepackage[left=26mm,right=26mm,top=24mm,bottom=24mm,headheight=14pt]{geometry}
\usepackage{xcolor,fancyhdr}
\usepackage[hypertexnames=false]{hyperref}
\definecolor{ink}{HTML}{1F4B60}
\hypersetup{colorlinks=true,linkcolor=ink,citecolor=ink,urlcolor=ink,
pdftitle={Mean, median and mode of stable laws, and a sharp gamma-mixture criterion},
pdfauthor={Research manuscript prepared with ChatGPT for Keio}}
\pagestyle{fancy}\fancyhf{}\lhead{\small\itshape Medians and modes of stable laws}
\rhead{\small\thepage}\renewcommand{\headrulewidth}{0.25pt}
\newtheorem{theorem}{Theorem}[section]
\newtheorem{lemma}[theorem]{Lemma}
\newtheorem{proposition}[theorem]{Proposition}
\newtheorem{corollary}[theorem]{Corollary}
\numberwithin{equation}{section}
\renewcommand{\Re}{\operatorname{Re}}\renewcommand{\Im}{\operatorname{Im}}
\setlength{\emergencystretch}{1.5em}
\title{Mean, median and mode of stable laws,\\
and a sharp gamma-mixture criterion}
\author{Research manuscript prepared with ChatGPT for Keio}
\date{3 November 2026}
\begin{document}
\maketitle
\begin{abstract}
For every asymmetric strictly stable law of index $1<\alpha<2$, classical
or free, the median lies strictly between the mean and the mode.
In the positivity parametrization a single explicit point separates
the median from the mode in both theories. The free proof evaluates
the tail along the boundary of the inverse Cauchy transform.
The classical proof tests the distribution function and the density slope
at the same point by Fourier inversion. One uniform lower bound is obtained
with a finite certificate in exact rational arithmetic, reproduced in full.
We also determine when the shifted positive free stable variable of index
$0<\alpha<1$ is a gamma scale mixture of shape $3/2$. The exact range is
$\alpha\le4/5$. Above this value an inverse-branch singularity contradicts
complete monotonicity. These results answer two questions of Hasebe, Simon
and Wang.
\end{abstract}

\section{Two questions on stable laws}
\label{sec:intro}
A skewed density can have different values for its mean, median and
mode. For stable densities, even deciding the order of these three
numbers has proved less straightforward than locating the mean.
Simon explicitly posed the classical question in 2014 \cite[Remark~12(b)]{Frechet}. Hasebe, Simon and Wang asked for the strict ordering when
$1<\alpha<2$ in both classical and free probability
\cite[Remark~14(b), p.~322]{HSW}. They pointed out that at maximum
positivity the mean is zero and the other two quantities are positive,
leaving the comparison between the median and mode. Nolan had
recorded numerical evidence for the classical ordering
\cite[Chapter~1, p.~14]{Nolan}. Biane's parametrization and unimodality
theorem describe the free mode \cite[Appendix]{BP}. The angular
parameter used below is that of Hasebe and Kuznetsov \cite{HK}.

Let $X$ denote the classical or the free strictly stable variable in
the normalizations specified in Sections~\ref{sec:free} and
\ref{sec:classical}. Write $\rho=\Pr(X>0)$, let $m$ be its median, and
let $M$ be the unique mode of its density. The range of $\rho$ is
$[1-1/\alpha,1/\alpha]$. Consider the auxiliary problem of finding
one explicit point that separates $m$ and $M$ for both laws with the
same parameters. Put
\begin{equation}\label{eq:separator}
 q=\pi\alpha(\rho-1/2),\qquad
 c=(\cos q)^{1/\alpha}\tan q .
\end{equation}
\begin{theorem}[Ordering for classical and free stable laws]\label{thm:order}
For $1<\alpha<2$ the following relations hold, including at both
extremal skewness parameters:
\[
\begin{array}{ccl}
\rho>1/2 &\implies& \mathbb E X=0<m<c<M,\\
\rho<1/2 &\implies& M<c<m<0=\mathbb E X,\\
\rho=1/2 &\implies& \mathbb E X=m=c=M=0 .
\end{array}
\]
For a location and positive scale change, the same change applied to
the four quantities preserves their order.
\end{theorem}
At $\alpha=2$ the classical Gaussian and free semicircle laws are symmetric,
so their means, medians and modes coincide.
For example, if $\alpha=3/2$ and $\rho=2/3$, both normalized laws have
median in $(0,2^{-1/3})$ and mode strictly above $2^{-1/3}$.
In the classical family $c$ is the location obtained by converting
the centered strictly stable parametrization to Nolan's $S_0$
parametrization \cite[Chapter~1, \S1.3]{Nolan}. The comparison with
the median and mode gives this familiar location another interpretation.

Our second result concerns the free one-sided law $X_\alpha=X_{\alpha,1}$
with $0<\alpha<1$. Its lower support endpoint is
$b_\alpha=\alpha(1-\alpha)^{1/\alpha-1}$. Let $V_d$ have the gamma
density $x^{d-1}e^{-x}/\Gamma(d)$ on $x>0$.
\begin{theorem}[A sharp mixture threshold]\label{thm:mixture}
There is a positive random variable $S$ independent of $V_{3/2}$ with
$X_\alpha-b_\alpha\stackrel{d}{=}S V_{3/2}$ if and only if
$0<\alpha\le4/5$.
\end{theorem}
Hasebe, Simon and Wang proposed this representation for every
$0<\alpha<1$ and established it for $\alpha\le3/4$
\cite[Remarks~2(b) and 10(b)]{HSW}. The positive direction through
$4/5$ follows from Wang's Thorin classification \cite[Theorem~C.4]{MW}
and the square-root density at the edge. Our working manuscript \emph{Restricted partitions, domination polynomials,
and the Kanter factor} (3 November 2026) gave a coefficient obstruction for
$23/25\le\alpha<1$. Here we close the interval $(4/5,23/25)$ and
obtain an oscillation statement for derivatives. Wang also excludes
the generalized gamma convolution property above $4/5$
\cite[Theorem~A.3]{MW}. A fixed-shape gamma scale mixture need not
be a generalized gamma convolution, so that exclusion alone does not
prove the necessity in Theorem~\ref{thm:mixture}.
\section{The free stable law}\label{sec:free}
We first prove the free part of Theorem~\ref{thm:order}. The argument
also gives the tail probability at the mode.
\begin{proof}[Proof of the free case]
Write $X=X_{\alpha,\rho}$, $p=\alpha-1$, $A=\pi\alpha\rho$, and $F=1/G$ for the reciprocal Cauchy transform, with $G(z)=\mathbb E[(z-X)^{-1}]$. The inverse transform is $z=H(F)=F-e^{iA}F^{-p}$, with the power on the upper half-plane. On the interior of the support, write
\[
 F=R e^{i\theta},\qquad v=A-p\theta,\qquad 0<\theta<\pi .
\]
The boundary parametrization of Biane \cite[Appendix~A]{BP} also follows directly by setting $\operatorname{Im} H=0$. Since $p\pi\le A\le\pi$ and $0<p<1$, we have $0<v<\pi$ and obtain
\[
 R^\alpha=\frac{\sin v}{\sin\theta},\qquad
 x(\theta)=R\frac{\sin(v-\theta)}{\sin v},\qquad
 f(x(\theta))=\frac{\sin\theta}{\pi R}.
\]
The last equality is Stieltjes inversion. In the coordinates $(\log R,\theta)$ the imaginary part of $H$ has radial derivative $\alpha R\sin\theta$ along the curve. The Cauchy--Riemann equations therefore give
\[
 x'(\theta)=-\frac{|F H'(F)|^2}{\alpha R\sin\theta}<0.
\]
The numerator is positive because $\operatorname{Im}(FH'(F))=\alpha R\sin\theta$. At $\theta\downarrow0$, the parameter $x$ tends to $+\infty$ when $A<\pi$, and to the finite right edge $\alpha p^{1/\alpha-1}$ when $A=\pi$. At $\theta\uparrow\pi$, the limit is $-\infty$ when $A>p\pi$, and is the reflected finite edge when $A=p\pi$. There are no atoms and the density is positive on the interior of this interval.

Along this boundary an antiderivative of $G\,dz$ is
\[
 J=\log F-\frac p\alpha e^{iA}F^{-\alpha}.
\]
Its imaginary part tends to zero at the right endpoint. When the right endpoint is finite, $F$ there is real and $A=\pi$. At an infinite right endpoint $R\to\infty$ and $\theta\to0$. Integrating $-\operatorname{Im}G/\pi$ from $x(\theta)$ to the endpoint gives
\begin{equation}\label{eq:free-median-tail}
 \Pr\{X>x(\theta)\}=\frac1\pi\left[
 \theta-\frac p\alpha\frac{\sin\theta}{\sin v}\sin(v-\theta)\right].
\end{equation}
At $\theta=\pi\rho$, $v=\theta$, so $x=0$ and the tail is $\rho$.
At $\theta=\pi/2$, $v=\pi/2+q$, so $x=c$ and
\[
 \Pr\{X>c\}=\frac12-\frac{p}{\pi\alpha}\tan q .
\]

For completeness, the density's unique mode follows directly by logarithmic differentiation. Its angular derivative has the sign of
$S(\theta)=(\alpha+1)\cot\theta+p\cot v$.
As $v$ lies between $p(\pi-\theta)$ and $\pi-p\theta$,
concavity of sine on $[0,\pi]$ gives
$\sin v\ge\min\{\sin(p(\pi-\theta)),\sin(p\theta)\}\ge p\sin\theta$.
Consequently
\[
 S'=-(\alpha+1)\csc^2\theta+p^2\csc^2 v\le-\alpha\csc^2\theta<0.
\]
The density tends to zero at both ends, so it has one interior maximum. If $\rho>1/2$, then $q>0$ and $S(\pi/2)=-p\tan q<0$. Because $x$ decreases with $\theta$, $M>c$. The two tail evaluations already give $0<m<c$. As a further identity, if $t$ is the unique root of $S$, then $p\cot v=-(\alpha+1)\cot t$ there. Equation~\eqref{eq:free-median-tail} simplifies to
\[
 \Pr\{X>M\}=\frac{t-\sin(2t)}{\pi}.
\]
Reflection $X_{\alpha,1-\rho}\stackrel{d}{=}-X_{\alpha,\rho}$ handles $\rho<1/2$. Symmetry handles $\rho=1/2$.

Finally $X$ is integrable. On an unbounded right tail $\sin A>0$, $R\sim x$, and $\sin\theta=\sin v/R^\alpha$, hence $f(x)\sim (\sin A)/(\pi x^{\alpha+1})$. The other unbounded tail follows by reflection. The inverse-transform equation and $F(iy)\sim iy$ give $F(iy)-iy=O(y^{-p})$, hence $(iy)^2(G(iy)-1/(iy))\to0$. Dominated convergence identifies this limit with $\mathbb E X$, since $|iy X/(iy-X)|\le |X|$. This proves the claimed mean.
\end{proof}

\section{The classical stable law}\label{sec:classical}
In the corresponding classical normalization \cite{HK}, the
characteristic function of $Y$ for $t>0$ is
$\exp(-e^{-iq}t^\alpha)$, with $q$ from \eqref{eq:separator}.
It is given by conjugation for $t<0$.
\begin{proof}[Proof of the classical case]
For the classical law the absolute first moment is finite by the identity
$\mathbb E|Y|=(2/\pi)\int_0^\infty(1-\Re\mathbb E e^{itY})t^{-2}\,dt$,
obtained by Tonelli from the elementary cosine integral. The integrand is $O(t^{\alpha-2})$ at zero. Differentiating the characteristic function at zero now gives mean zero. We need only treat $q>0$.
After rescaling $Z=Y/(\cos q)^{1/\alpha}$ its characteristic function on $t>0$ is $\exp[-t^\alpha+iB t^\alpha]$, where $B=\tan q>0$. Denote its density by $f$. Set $p=\alpha-1$ and $D=pB$. Admissibility of $\rho$ gives
\[
 D\le p\cot(\pi p/2)\le2/\pi<2/3.
\]
The known unimodality of stable densities \cite{Simon} and positivity of $f$ allow us to use a local derivative to locate the mode. Indeed the density is analytic by exponential decay of the characteristic function, so a modal plateau is impossible.

Fourier inversion gives the density derivative, and inversion for the continuous distribution gives the tail. At the proposed separator $x=B$ these are
\[
 \Pr(Z>B)-\tfrac12=(D/\pi)I_{-1},\qquad f'(B)=(D/\pi)I_1,
\]
where
\[
 I_j=\int_0^\infty t^j e^{-t^{1+p}}h_p(t)\,\operatorname{sinc}(Dh_p(t))\,dt,
 \qquad h_p(t)=t(t^p-1)/p,\qquad \operatorname{sinc}(u)=\sin(u)/u.
\]
For verification only, extend $h_p$ to $p=0$ by $h_0(t)=t\log t$. The signs will hold throughout the larger parameter rectangle $0<p\le1$, $0<D\le2/3$.

On $0<t<1$ we have $|h_p(t)|\le-t\log t\le1/e<1/2$. Consequently $\operatorname{sinc}(Dh_p)\ge1-1/54$ there. Since $h_p\ge0$ on $t\ge1$, $\operatorname{sinc}\le1$, and $\int_0^1 |h_p|/t\,dt\le1$, we have
\[
 I_{-1}\le \frac{1-\Gamma(1/\alpha)}{\alpha p}+\frac1{54}
 \le-\frac{\gamma}{\alpha^2}+\frac1{54}<0 .
\]
The first term is the integral with the sinc factor replaced by $1$. The second inequality is convexity of $\Gamma$ and its tangent at $1$, with $\Gamma'(1)=-\gamma$. The elementary bound $\gamma\ge1-\log 2>1/4$ suffices.

For $I_1$, on $t\le1$ use $e^{-t^{1+p}}\le e^{-t^2}\le 1-t^2/2$ and $h_p\le0$. This part is at least
\[
 -\frac1{3(3+p)}+\frac1{10(5+p)}.
\]
For $h\ge0$, a lower bound on $\operatorname{sinc}(Dh)$ valid uniformly in $D$ is
\[
 L(h)=\begin{cases}
 \max(0,1-2h^2/27),&0\le h\le9/2,\\
 -1/4,&h>9/2 .
 \end{cases}
\]
In the first range, use $Dh\le3<\pi$ and $\sin v\ge v-v^3/6$ for $v\ge0$. In the second, $\operatorname{sinc}(v)\ge-1/4$ holds for all $v\ge0$. On $[\pi,4]$ use $\sin v\ge-(v-\pi)\ge-v/4$, and for $v\ge4$ use $\sin v\ge-1$. Here $v=0$ is interpreted by continuity. For $t\ge1$, $h_p(t)$ increases with both $p$ and $t$, since $h_p=t\log t\int_0^1 t^{pu}\,du$. In particular $h_p\le h_1=t(t-1)$. The integral beyond $T=24$ is accordingly bounded below by
$-e^{-T}(T^3+3T^2+6T+6)/4$, by integration by parts.

It remains to bound the integral from $1$ to $24$. Partition the $p$-interval into 32 equal intervals $[l,u]$ and the $t$-interval into intervals $[a,b]$ of width $1/64$. Put $H=h_u(b)$. When $H\le9/2$, a lower bound for the cell integral is
\[
 (b-a)\,a\,e^{-b^{1+u}}h_l(a)\max(0,1-2H^2/27).
\]
Otherwise it is $-(b-a)b e^{-a^{1+l}}H/4$. These inequalities follow by monotonicity of $h$, of $t^{1+p}$ on this domain, and of $L$. Sum these certified lower bounds, together with the two endpoint bounds, using $-1/(3(3+l))+1/(10(5+u))$ at the lower endpoint. The exact calculation in Appendix~\ref{sec:certificate} gives a total exceeding $1/100$ for each of the 32 parameter intervals. This rigorously proves $I_1>0$.

Thus the median of $Z$ is below $B$ and the mode is above $B$. Positivity of the median follows from $\Pr(Z>0)=\rho$; directly, the same tail inversion at zero gives $1/2+(\pi\alpha)^{-1}\int_0^\infty e^{-s}\sin(Bs)\,ds/s=1/2+\arctan(B)/(\pi\alpha)=\rho$. Undoing the scale proves the theorem, including reflection and symmetry.
\end{proof}
\section{The obstruction to gamma mixing above four fifths}
\label{sec:mixture-obstruction}

Write $X_\alpha=X_{\alpha,1}$ in the parametrization of
Hasebe, Simon and Wang \cite{HSW}, let $f_\alpha$ denote its density,
and put
\[
 p=\frac{\alpha}{1-\alpha},\qquad m=p+1=\frac{1}{1-\alpha},
 \qquad P(t)=t-t^{m},\qquad t_0=m^{-1/p},\qquad
 b_\alpha=P(t_0)=\frac{p}{m}t_0 .
\]
Complex powers of $t$ below use $-\pi<\arg t<\pi$.
The point $b_\alpha=\alpha(1-\alpha)^{1/\alpha-1}$ is the lower
support endpoint of $X_\alpha$. In this section write $b=b_\alpha$.

\begin{proposition}\label{prop:phase-nonmixture}
If $4/5<\alpha<1$, the function
\[
 x\longmapsto \frac{f_\alpha(b+x)}{\sqrt x},\qquad x>0,
\]
is not completely monotone. Hence $X_\alpha-b$ is not a scale
mixture of a gamma law of shape $3/2$.
\end{proposition}

The assertion about scale mixtures follows from their
densities. If $S>0$ is independent of a unit-scale gamma variable of
shape $3/2$, the density of its product with $S$, divided by $\sqrt{x}$,
equals $\mathbb E[S^{-3/2}e^{-x/S}]/\Gamma(3/2)$.
We prove the proposition by exhibiting a singular continuation of
the normalized density inside a half-plane where complete
monotonicity would force a holomorphic continuation.

\subsection{Continuation on the positive real axis}
Let $G(z)=\mathbb E[(z-X_\alpha)^{-1}]$ and $F=1/G$. The
Voiculescu transform in \cite[equation~(2)]{HSW} gives
$z=F-e^{i\pi\alpha}F^{1-\alpha}$.
This holds first in the defining inverse domain and then throughout
the upper half-plane by analytic continuation, since $F$ takes
values there. At $z<b$ on the real axis, $G(z)<0$ and $F(z)<0$.
Set $t=(-F)^{1/m}>0$. Using the upper boundary value of $F^{1-\alpha}$
yields the real equations
\begin{equation}\label{eq:mixture-inverse}
 z=P(t)=t-t^m,\qquad G(z)=-t^{-m}.
\end{equation}
The function $P$ increases on $(0,t_0)$ and decreases on
$(t_0,\infty)$. For $0<z<b$ write $t_{\rm in}(z)<t_0<t_{\rm out}(z)$
for its two positive inverses. The reciprocal Cauchy transform
selects $t_{\rm out}$ in \eqref{eq:mixture-inverse}. For $z<0$
the positive solution is unique and exceeds $1$, and this solution
continues through $t=1$ at $z=0$.

At $t_0$ we have $P'(t_0)=0$ and $P''(t_0)=-mp\,t_0^{p-1}<0$.
Thus $b-P(t)=(t-t_0)^2 B(t)$, with $B$ real analytic and
$B(t_0)>0$. Inverting the coordinate $(t-t_0)\sqrt{B(t)}$ gives
a convergent expansion $t=T(\xi)$ around $\xi=0$, where
$\xi^2=b-z$, $T(0)=t_0$, $T'(0)>0$, and
$t_{\rm out}(z)=T(\sqrt{b-z})$ on the real interval immediately
to the left of $b$. Set $Q(t)=-t^{-m}$. The quotient
\[
 \frac{Q(T(\xi))-Q(T(-\xi))}{2\pi\xi}
\]
is analytic and even in $\xi$. At $z=b+x+i0$, $x>0$ small, the
physical root has $\xi=-i\sqrt x$. The two values in the numerator
are conjugates, and by Stieltjes inversion the quotient is
$f_\alpha(b+x)/\sqrt x$. In particular this normalized density
extends holomorphically through the edge. To its left the
continuation equals
\begin{equation}\label{eq:mixture-branch-difference}
 H(z)=\frac{Q(t_{\rm out}(z))-Q(t_{\rm in}(z))}
                {2\pi\sqrt{b-z}},\qquad 0<z<b .
\end{equation}
Both inverses there are regular. On $z>b$ the density is real
analytic, as recalled in \cite[Section~1]{HSW}. We have obtained a
single real-analytic function $H$ on $(0,\infty)$ whose values for
$z>b$ are $f_\alpha(z)/\sqrt{z-b}$.

\subsection{A half-plane forced by the Taylor signs}
Suppose $x\mapsto H(b+x)$ were completely monotone.
Fix $c>b$ and develop $H(c-w)$ at $w=0$. Its coefficients
$(-1)^n H^{(n)}(c)/n!$ are nonnegative. The germ has a real-analytic
continuation along $0\le w<c$, by the preceding construction.
Consequently its Taylor radius is at least $c$. Indeed, if the radius
were $R<c$, the continuation would be holomorphic in a neighborhood
of the positive point $R$. For real $r<R$ sufficiently close to $R$,
the Taylor series at $r$ would then converge at some $r+s>R$,
$s>0$. Expanding $(r+s)^n$ and interchanging nonnegative sums shows
that the original Taylor series must converge there as well. This
is the positive-coefficient singularity principle of Pringsheim.

It follows that the germ at $c$ extends to $|z-c|<c$. These disks
are nested as $c$ increases and their extensions agree on overlaps
by the identity theorem, since they agree with $H$ on the common
positive real interval. Their union is the entire right half-plane. For $\Re z>0$, choose $c>\max\{b, |z|^2/(2\Re z)\}$. Complete
monotonicity would therefore give a single holomorphic extension of
$H$ to $\Re z>0$.

\subsection{A second critical value}
Assume $p>4$ and put $\varphi=2\pi/p$ and $\omega=e^{i\varphi}$.
Then $0<\varphi<\pi/2$ and
\[
 z_*=b\omega\quad\hbox{satisfies}\quad \Re z_*>0 .
\]
Continue \eqref{eq:mixture-branch-difference} from $0<z<b$ toward
a small positive $\varepsilon$, around the circular arc
$\varepsilon e^{i\eta}$, $0\le\eta\le\varphi$, and then on the ray
$z=s\omega$, $\varepsilon\le s<b$. This entire path lies in
$\Re z>0$ and avoids $b$. We determine both inverse branches on
the last segment.

Near zero the small inverse has the form
$t_{\rm in}(z)=z A(z^p)$, with $A(0)=1$. Indeed $A=1+\zeta A^m$
has a unique holomorphic solution at $\zeta=0$ with this value,
by the implicit function theorem. The powers in $z^p$ can be
continued throughout the small circular arc. Since $p\varphi=2\pi$,
at its far end the value of $t_{\rm in}$ is
$\omega t_{\rm in}(\varepsilon)$. For all $0<s<b$ one has
$P(\omega t_{\rm in}(s))=\omega P(t_{\rm in}(s))=\omega s$;
the arguments of these $t$-values are $\varphi$. Local uniqueness
now identifies the continued small branch as $\omega t_{\rm in}(s)$
on the full ray. As $s\uparrow b$ it reaches $\omega t_0$ with a
square-root singularity whose leading coefficient is nonzero.
The derivative $Q'(\omega t_0)=m(\omega t_0)^{-m-1}$ is nonzero,
so composition with $Q$ retains this singularity.

The other inverse stays regular at $z_*$. To see this, put
$v=1-t^p$. Then $z=t v=v(1-v)^{1/p}$, so the solution near $t=1$
is $t_{\rm out}=(1-v)^{1/p}$, with
\[
                 v=z(1-v)^{-1/p},\qquad v(0)=0.
\]
Lagrange inversion gives
$[z^n]v(z)=n^{-1}[u^{n-1}](1-u)^{-n/p}>0$ for every $n\ge1$.
On $0<z<b$ the real inverse increases from zero to $p/(p+1)$
as $z\uparrow b$. On this interval its derivative is finite because
$v(1-v)^{1/p}$ has nonzero derivative for $0\le v<p/(p+1)$.
The positive-coefficient argument just proved shows that the radius
of the Taylor series for $v$ is at least $b$. Monotone convergence now
gives $\sum_{n\ge1}([z^n]v)b^n=p/(p+1)$, so the series converges
absolutely and uniformly on the closed disk. At $z_*=b\omega$,
strict triangle inequality yields $|v(z_*)|<p/(p+1)$ because the
first two nonzero summands have different arguments. In particular
$v(z_*)\ne p/(p+1)$ and $|v(z_*)|<1$. The equation for $v$ passes
to the boundary by continuity, and the inverse function theorem
extends $v$ holomorphically there. Throughout the disk $1-v$ stays
in the right half-plane. The same choice of its $1/p$ power therefore
continues the outer inverse and $Q(t_{\rm out})$ along the path and
through $z_*$. 

In \eqref{eq:mixture-branch-difference} precisely one numerator term
has a nonzero square-root term at $z_*$. The denominator has a
holomorphic, nonzero local branch there because $z_*\ne b$.
The resulting continuation of $H$ cannot be holomorphic at this
point. This contradicts the half-plane extension and proves
Proposition~\ref{prop:phase-nonmixture}.

\begin{corollary}\label{cor:oscillate}
Let $4/5<\alpha<1$ and
$c>\max\{b,b/(2\cos(2\pi/p))\}$. Then
$(-1)^n H^{(n)}(c)$ takes strictly positive values for infinitely many
$n$ and strictly negative values for infinitely many $n$.
\end{corollary}
\begin{proof}
The path in the preceding argument can be taken in $|z-c|<c$:
on the ray $z=s\omega$ this is exactly $s<2c\cos(2\pi/p)$,
and all angles of the small initial arc have cosine at least
$\cos(2\pi/p)>0$. The entire real interval $(0,b]$ also lies in
the disk. If either sign were eventually absent, deletion of finitely
many Taylor terms from $H(c-w)$, and a change of sign if necessary,
would give nonnegative coefficients. The same Taylor-radius argument
forces holomorphy on the disk, contradicting the singularity at $z_*$.
\end{proof}
For example, $\alpha=5/6$ gives $p=5$. One may take $c=2b$ since
$\cos(2\pi/5)=(\sqrt5-1)/4>1/4$. Thus infinitely many derivative
tests fail already at the single point $x=b$ of the shifted density
divided by $\sqrt{x}$.
\section{The scale mixture up to the threshold}
\label{sec:mixture-positive}

Suppose $0<\alpha\le4/5$. In his October 2026 preprint Wang
\cite[Theorem~1.3 and Theorem~C.4]{MW} proves that
$X_{\alpha,1}$ belongs to the extended Thorin class. We
derive the needed gamma shape from the behavior of its density at the
lower endpoint. This also explains why merely having a positive
Thorin representation does not specify the shape of the mixture.

Let $Y=X_\alpha-b_\alpha$ and $L(s)=\mathbb E[e^{-sY}]$.
Translation preserves the extended Thorin property. The
L\'evy--Khintchine representation for a nonnegative infinitely
divisible variable gives a nonnegative drift and positive jumps.
Its L\'evy density in the Thorin class is $k(x)/x$ for a completely
monotone $k$. By Bernstein's theorem there is a positive measure
$U$ with $k(x)=\int e^{-rx}U(dr)$. It has no atom at $r=0$, since
otherwise $\int_1^\infty k(x)\,dx/x$ would diverge. Tonelli's
theorem and Frullani's identity give
\begin{equation}\label{eq:mixture-thorin}
 -\log L(s)=ds+\int_{(0,\infty)}\log(1+s/r)\,U(dr).
\end{equation}
To see that the drift is zero, note that for every $\varepsilon>0$ the event
$Y<\varepsilon$ has positive probability, so
$0\le d\le\limsup_{s\to\infty}(-\log L(s))/s\le\varepsilon$.

The local inverse at the edge derived in
Section~\ref{sec:mixture-obstruction} works for every $0<\alpha<1$.
As $y\downarrow0$ it gives the density of $Y$ as
$c\sqrt y+O(y^{3/2})$ with $c>0$. In fact the leading coefficient is
$Q'(t_0)T'(0)/\pi>0$ in the notation of that section. Laplace
integration of the leading powers gives
\[
 L(s)\sim c\,\Gamma(3/2)s^{-3/2},
 \qquad -L'(s)\sim c\,\Gamma(5/2)s^{-5/2}
 \qquad(s\to\infty).
\]
For these two estimates it suffices to split the density integral at
a fixed small positive number. The complementary integral, also
after multiplication by $y$, is exponentially small. Differentiation
of \eqref{eq:mixture-thorin} for $s>0$ gives
\[
 -s\frac{L'(s)}{L(s)}
 =\int_{(0,\infty)}\frac{s}{s+r}\,U(dr)
 \ \uparrow\ U((0,\infty))
 \qquad (s\to\infty).
\]
For differentiation locally in $s>0$ use
$s/(s+r)\le\log(1+s/r)$ and the finiteness in
\eqref{eq:mixture-thorin}. The asymptotics prove that the total
Thorin mass is $D=\Gamma(5/2)/\Gamma(3/2)=3/2$.

For completeness, finite Thorin mass $D$ yields a shape-$D$ scale
mixture as follows. First suppose $U=\sum_i d_i\delta_{r_i}$ has
finite support. The variable in \eqref{eq:mixture-thorin} with zero
drift has the same law as $\sum_i V_i/r_i$, where $V_i$ are independent
unit-scale gamma variables of shapes $d_i$. The change of variables
$V_i=V P_i$, $V=\sum_i V_i$, shows that $V$ is gamma of shape
$D=\sum_i d_i$ and independent of $(P_i)_i$. Thus the mixing scale
is $\sum_i P_i/r_i$.

For general $U$ of finite mass $D$, push it forward under
$r\mapsto\min(n,\max(1/n,\lceil nr\rceil/n))$. The resulting
finite-support measures have total mass $D$, and their transforms
in \eqref{eq:mixture-thorin} converge pointwise by dominated
convergence. Indeed the new integrand is bounded by
$\log(1+s/r)+\log(1+s)$, an integrable majorant. The corresponding
laws $V S_n$ converge weakly to $Y$ by the continuity theorem for
Laplace transforms, with $V$ unit-scale gamma of shape $D$ independent
of $S_n$. The scales are tight because
\[
 \mathbb P(S_n>R)\,\mathbb P(V\ge1)\le\mathbb P(VS_n>R).
\]
Pass to a weakly convergent subsequence of $S_n$. Its limit $S$,
taken independent of $V$, gives $Y\stackrel{d}{=}VS$. Since $Y$
has a density, $S>0$ almost surely. This proves the scale-mixture
assertion on the closed parameter range $0<\alpha\le4/5$.

\paragraph{Connection with the source corpus.}
Manuscript 002 of the supplied corpus studies continuation of inverse
branches around critical values \cite{corpus}. In the present
application one has to distinguish the two branches entering a
\emph{density}: the singularity excludes a gamma mixture only after
one verifies that it survives their subtraction. The proof makes that
test explicitly for the free stable inverse $t-t^{1/(1-\alpha)}$.
\appendix
\section{Exact lower bounds for the classical Fourier integral}
\label{sec:certificate}
The verifier below is complete Python~3 code and uses only its
standard library. It encloses real powers and exponentials by rational
numbers. For $S=2^{60}$, five successive integer square roots give
$\lfloor S t^{j/32}\rfloor$ from $\lfloor S^{32}t^j\rfloor$.
For $N=2^{18}$ and $0\le x<N$, the two sides of
\[
              (1-x/N)^N\le e^{-x}\le(1+x/N)^{-N}
\]
are evaluated by squaring with directed rounding. At $p=0$ the
positive cell bound uses $\log t\ge 2(v+v^3/3)$, where
$v=(t-1)/(t+1)\in[0,1)$, by the power series for $\operatorname{arctanh}v$.
The code sums lower bounds over 1,472 integration cells for each of
32 parameter intervals. The smallest resulting lower bound, rounded
down to six decimal places, is $0.017207$, on $31/32\le p\le1$.
Only the comparison with $1/100$ is used.
\begin{Verbatim}[fontsize=\small]
from fractions import Fraction as Q
from math import isqrt

S = 2**60

def floor(x):
    return x.numerator // x.denominator

def ceil(x):
    return -floor(-x)

def power_bounds(t, j):
    # Dyadic bounds on t**(j/32) by five integer square roots.
    n = floor(t**j * S**32)
    for _ in range(5):
        n = isqrt(n)
    return Q(n, S), Q(n+1, S)

def exp_minus(x, upper=False):
    # (1-x/N)**N <= exp(-x) <= (1+x/N)**(-N).
    steps = 18
    u = x / (2**steps)
    assert 0 <= u < 1
    n = ceil(S/(1+u)) if upper else floor(S*(1-u))
    for _ in range(steps):
        n = (n*n + (S-1 if upper else 0)) // S
    return Q(n, S)

def h_lower(t, j, r):
    if j:
        return t*(r-1)*32/j
    v = (t-1)/(t+1)
    return 2*t*(v + v**3/3)

T, K = 24, 64
tail = exp_minus(Q(T), True)*(T**3+3*T**2+6*T+6)/4
bounds = []
for j in range(32):
    p, q = Q(j,32), Q(j+1,32)
    low = -1/(3*(3+p)) + 1/(10*(5+q)) - tail
    ticks = 0
    for i in range((T-1)*K):
        a, b = 1+Q(i,K), 1+Q(i+1,K)
        r = power_bounds(a,j)[0]
        s = power_bounds(b,j+1)[1]
        hi = b*(s-1)/q
        if 2*hi <= 9:
            w = max(0, 1-2*hi**2/27)
            term = a*exp_minus(b*s)*h_lower(a,j,r)*w
        else:
            term = -b*exp_minus(a*r, True)*hi/4
        ticks += floor(S*term)
    low += Q(ticks,S*K)
    assert low > Q(1,100), j
    bounds.append(floor(10**6*low))
print(min(bounds), bounds.index(min(bounds)))
\end{Verbatim}
\begin{thebibliography}{9}\raggedright
\bibitem{BP} H.~Bercovici and V.~Pata, with an appendix by P.~Biane,
Stable laws and domains of attraction in free probability theory,
\emph{Ann. of Math.} (2) \textbf{149} (1999), 1023--1060.
\href{https://doi.org/10.2307/121080}{doi:10.2307/121080}.
\bibitem{HK} T.~Hasebe and A.~Kuznetsov, On free stable distributions,
\emph{Electron. Commun. Probab.} \textbf{19} (2014), no.~56, 1--12.
\href{https://doi.org/10.1214/ECP.v19-3443}{doi:10.1214/ECP.v19-3443}.
\bibitem{HSW} T.~Hasebe, T.~Simon and M.~Wang, Some properties of the
free stable distributions, \emph{Ann. Inst. Henri Poincar\'e Probab.
Statist.} \textbf{56} (2020), 296--325. First posted in 2018.
\href{https://doi.org/10.1214/19-AIHP962}{doi:10.1214/19-AIHP962}.
\bibitem{Nolan} J.~P.~Nolan, \emph{Stable Distributions. Models for
Heavy Tailed Data}, Chapter~1, 2018 manuscript.
\href{https://edspace.american.edu/jpnolan/wp-content/uploads/sites/1720/2020/09/Chap1.pdf}{Author's text}.
\bibitem{Frechet} T.~Simon, Comparing Fr\'echet and positive stable laws,
\emph{Electron. J. Probab.} \textbf{19} (2014), no.~16, 1--25.
\href{https://doi.org/10.1214/EJP.v19-3058}{doi:10.1214/EJP.v19-3058}.
\bibitem{Simon} T.~Simon, A multiplicative short proof for the unimodality
of stable densities, \emph{Electron. Commun. Probab.} \textbf{16}
(2011), 623--629.
\href{https://doi.org/10.1214/ECP.v16-1672}{doi:10.1214/ECP.v16-1672}.
\bibitem{MW} M.~Wang, \emph{Classical infinite divisibility,
self-decomposability and bell-shape of free stable laws},
preprint, 9 October 2026.
\href{https://hexagonmath.org/pdf/2610.00169v1}{Hexagon 2610.00169v1}.
\bibitem{corpus} QEEDDD, \emph{Verifiable Mathematical Reasoning
Corpus}, manuscript 002, \emph{The Polyexponential Trace-Tower Program},
revision \href{https://github.com/QEEDDD/Verifiable-Mathematical-Reasoning-Corpus/tree/ecd87de3348dcb706a76bb12fde722d7a348a012}{\texttt{ecd87de3348dc}}.
\end{thebibliography}
\end{document}
